\lesson{II.6}{How fast should you believe?}{An observer that believes its sensor quickly catches every push — and every speck of noise. Where the balance lies, and what finally breaks it.}{1}{adrc-bw}{Part II · Beyond PID}
\label{les:adrcbw}

\lead{\setupcard{\setuprow{Level}{L1}\setuprow{Controller}{ADRC, $\omega_c = 3$, with a fast observer, $\omega_o = \SI{40}{rad/s}$}\setuprow{Altimeter}{\SI{2}{cm} of noise}\setuprow{Wind}{on, as in Lesson~1}\setuprow{Goal}{a score below 7.5: RMS error in cm plus ten times the thrust jitter in N}}%
The observer of the last lesson had one knob, its bandwidth $\omega_o$. Turned up, it sees a disturbance sooner. But the observer cannot tell a push on the drone from an error of the altimeter: both show up as a surprise in the measurement. This lesson turns the knob from 4 to 70 radians per second and measures both sides of the trade. The score has a clear minimum, the noise side has a formula, and beyond the useful range lie two different kinds of failure.}

\section{Two kinds of surprise}

The observer of \cref{eq:adrc-eso} corrects its estimate with $L\,(y - \hat y)$. The difference $y - \hat y$, the innovation, has two sources: the drone moved in a way the model did not predict — a disturbance — or the sensor reported something the drone did not do — noise. The observer treats both the same way.\pitfall{Turning up an observer to catch faster disturbances also turns up every flaw of its sensor. A sharp observer on a poor sensor is worse than a dull one.} With a large $\omega_o$ its gains are large, and it reacts strongly to every surprise.

The two effects have opposite trends.
\begin{itemize}
  \item \textbf{Disturbances} are noticed after a few $1/\omega_o$ (Example~\ref{ex:adrc-eso}). A faster observer leaves less time for a gust to push the drone: the altitude error falls.
  \item \textbf{Noise} passes from the altimeter into $\hat f$, and $\hat f$ goes straight into the thrust. A faster observer amplifies it more: the thrust jitters.
\end{itemize}
The lesson's score weighs the two: the RMS altitude error in centimetres plus ten times the thrust jitter in newtons, the jitter being the RMS change of the thrust between samples \SI{10}{ms} apart.\tip{A score that adds two unlike things is a statement of priorities. \enquote{Ten times} says that a newton of jitter is as bad as ten centimetres of error. Change the factor and the best $\omega_o$ moves.}

\simfig{adrcbw_score}{Twelve flights, one for each observer bandwidth, with \SI{2}{cm} of altimeter noise and the wind of Lesson~1. Left: the altitude error falls and the thrust jitter rises; the dashed line is the jitter predicted from the noise alone (Example~\ref{ex:adrcbw-jitter}). Right: the score is smallest at $\omega_o = 6$. Beyond about 35 the noise drives the motors into their limits and both halves grow.}{fig:adrcbw-score}

\section{The noise side, exactly}

How strongly does the observer amplify noise? Its estimate of the disturbance is, in effect, the measured altitude differentiated twice and filtered: $f$ is the part of $\ddot y$ that the thrust does not explain.

\begin{example}[From altimeter noise to $\hat f$]
\label{ex:adrcbw-path}
Find the transfer function from the altimeter's noise $n$ to $\hat f$, and its gain at the observer's bandwidth.

\textbf{Step 1 — the observer with noise.} The observer sees $y + n$. Its error equations of Example~\ref{ex:adrc-eso} get the input $n$ through $L$: in the transfer function from the measurement to $\hat f$ only the third row matters.

\textbf{Step 2 — the result.} Working through the companion form (or directly: $\hat f$ is $\omega_o^3/(s + \omega_o)^3$ times the true $f$, and $f$ seen through the sensor is $s^2$ times the measurement):
\begin{equation}\label{eq:adrcbw-G}
  \frac{\hat f}{n}(s) = \frac{\omega_o^3\,s^2}{(s + \omega_o)^3} .
\end{equation}

\textbf{Step 3 — its gain.} At low frequency it rises like $\omega^2$ — a double differentiator. At $\omega = \omega_o$ it is $\omega_o^5/(2\sqrt2\,\omega_o^3) = \omega_o^2/2.83$. Far above it falls like $\omega_o^3/\omega$. For $\omega_o = 10$ the peak gain is about $\SI{35}{(m/s^2)/m}$; for $\omega_o = 40$, \SI{570}{(m/s^2)/m}: a centimetre of noise at that frequency becomes \SI{5.7}{N} of thrust for the \SI{1}{kg} drone.
\end{example}

\begin{theorem}[How the noise grows with the bandwidth]
\label{thm:adrcbw-scaling}
Let white noise of intensity $q$ enter the measurement of the observer \eqref{eq:adrc-eso}. Then the variance of the noise in $\hat f$ is
\begin{equation}\label{eq:adrcbw-var}
  \sigma_{\hat f}^2 = \frac{3}{16}\,q\,\omega_o^5 , \qquad\text{so}\qquad \sigma_{\hat f} \propto \omega_o^{5/2} .
\end{equation}
\end{theorem}
\begin{proof}
By \cref{thm:E-psd}, $\sigma^2 = \frac{q}{2\pi}\int|G(j\omega)|^2\,\dd\omega$ with $G$ from \cref{eq:adrcbw-G}. Substitute $\omega = \omega_o x$: $|G|^2 = \omega_o^4\,x^4/(1 + x^2)^3$ and $\dd\omega = \omega_o\,\dd x$, so $\sigma^2 = \frac{q\,\omega_o^5}{2\pi}\int_{-\infty}^\infty \frac{x^4}{(1 + x^2)^3}\,\dd x = \frac{q\,\omega_o^5}{2\pi}\cdot\frac{3\pi}{8}$.
\end{proof}

Doubling the observer bandwidth multiplies the noise in the disturbance estimate by $2^{5/2} = 5.7$. Between $\omega_o = 10$ and 20 the simulator's jitter grows by a factor of 4.9, between 20 and 30 by 2.6, where $1.5^{5/2} = 2.8$: the formula's exponent, slightly softened by the sampling.\recall{Appendix~E: a sequence of independent samples of variance $\sigma^2$ every $\Delta t$ acts like continuous white noise of intensity $q = \sigma^2\Delta t$.}

\begin{example}[The jitter, sample by sample]
\label{ex:adrcbw-jitter}
Predict the thrust jitter of the simulator for each $\omega_o$, from the noise alone.

\textbf{Step 1 — the sampled loop.} The controller runs at \SI{250}{Hz}. Put the plant (with its motor lag, discretised exactly), the observer (discretised exactly, as in \src{src/estimation/eso.ts}) and the law into one linear recursion, as in Example~\ref{ex:slow-discrete}.

\textbf{Step 2 — one bad reading.} Feed it a single altimeter error of \SI{1}{m} and record the thrust command $h_k$ that follows: the impulse response from noise to thrust.

\textbf{Step 3 — the jitter.} The thrust is held for \SI{4}{ms}, the jitter is measured over \SI{10}{ms}, so it spans two or three controller samples. By \cref{thm:noise-variance}, the variance of $u_{k+j} - u_k$ is $\sigma^2\sum_i(h_{i+j} - h_i)^2$; average $j = 2$ and $j = 3$, with $\sigma = \SI{2}{cm}$.

\textbf{Step 4 — compare.}
\begin{center}\small\sffamily
\begin{tabular}{rcccccc}
\toprule
$\omega_o$ [rad/s] & 4 & 10 & 15 & 20 & 25 & 30\\
\midrule
predicted [N] & 0.057 & 0.372 & 0.925 & 1.81 & 3.09 & 4.80\\
simulator [N] & 0.058 & 0.371 & 0.923 & 1.81 & 3.09 & 4.76\\
\bottomrule
\end{tabular}
\end{center}
The thrust jitter is the noise, through a linear filter, and nothing else.

\textbf{Step 5 — where it stops.} At $\omega_o = 40$ the prediction is \SI{9.7}{N} and the simulator shows \SI{8.6}{N}. A jitter of that size swings the thrust from zero to the motor limit: the motors clip, the clipped noise has a non-zero mean (\lessonref{les:noise}, rectification), and the altitude error jumps from \SI{0.6}{cm} at $\omega_o = 30$ to \SI{4.4}{cm}, \SI{16}{cm} at 50 and \SI{55}{cm} at 70. The drone is not unstable; it is being shaken by its own observer.\innumbers[-36pt]{At $\omega_o = 70$ the thrust jitters by \SI{15}{N} — more than the drone's weight — for a sensor error of two centimetres.}
\end{example}

\section{The disturbance side, and the best compromise}

The altitude error falls as the observer speeds up: \SI{2.6}{cm} RMS at $\omega_o = 4$, \SI{1.1}{cm} at 10, \SI{0.6}{cm} between 25 and 30 (\cref{fig:adrcbw-score}, left). Two things limit it from below. Noise itself moves the drone a little: the impulse response predicts about \SI{0.4}{cm} of noise-induced error at $\omega_o = 10$. And the wind's fast fluctuations are beyond any observer that sees them only through the altitude.

\begin{example}[Finding the best bandwidth]
\label{ex:adrcbw-best}
Read the best $\omega_o$ off the measurements and check it against the rule of thumb.

\textbf{Step 1 — the scores.} $\omega_o = 4$: $2.62 + 0.58 = 3.2$. 6: $1.79 + 1.27 = 3.05$. 8: $1.35 + 2.29 = 3.6$. 10: $1.08 + 3.71 = 4.8$. 15: $0.76 + 9.2 = 10.0$. 40: $4.4 + 86 = 91$.

\textbf{Step 2 — the minimum.} At $\omega_o = 6$, with the goal of 7.5 met from 4 to about 12. The starting value of the lesson, 40, misses it twelvefold.

\textbf{Step 3 — the rule of thumb.} $\omega_o$ three to ten times $\omega_c = 3$ means 9 to 30. With \SI{2}{cm} of noise and this weighting, the best is at the low end of that range and slightly below it. With a better altimeter the noise term shrinks by \cref{eq:adrcbw-var} and the optimum moves up: the rule of thumb is a starting point, the sensor decides.\tip{Tune an observer on the real sensor, not on a simulated perfect one. The best bandwidth depends on the noise more than on anything else.}
\end{example}

\section{The other limit: sampling}

Remove the noise and the wind. How fast may the observer be then? The rule of thumb in the ADRC literature points at the actuator: an observer much faster than the motors, which it does not model, should eventually fail. In this simulator it fails for another reason.

\begin{example}[The edge in clean air]
\label{ex:adrcbw-edge}
Find the largest stable $\omega_o$ of the sampled loop without noise, with and without the motor lag.

\textbf{Step 1 — the loop as a matrix.} The recursion of Example~\ref{ex:adrcbw-jitter} is $\symbf z_{k+1} = M(\omega_o)\,\symbf z_k$ with a $7 \times 7$ matrix: three plant states, three observer states and the held command.

\textbf{Step 2 — the edge.} Its largest eigenvalue reaches magnitude one at $\omega_o = \SI{161}{rad/s}$. Removing the motor lag from the plant ($\tau_m \to 0$) moves the edge by less than \SI{0.1}{rad/s}.

\textbf{Step 3 — read it.} The lag does not limit this observer. ADRC feeds the observer the commanded thrust, and the lag, seen through it, is just a fast part of $f$, estimated and cancelled with the rest. What limits it is the sampling: at the edge $\omega_o\Delta t = 161 \cdot 0.004 = 0.65$ — the observer is trying to react within a fraction of one sample.

\textbf{Step 4 — check.} Flown in calm air with a step of \SI{0.5}{m}: $\omega_o = 60$ and 100 are clean, 150 overshoots by \SI{28}{cm} and rings, 250 never leaves the ground.

\textbf{Step 5 — the general rule.} A sampled observer needs $\omega_o\Delta t$ well below one; the sampled controllers of \lessonref{les:slow} had the same limit in the form \enquote{sample ten to twenty times faster than the bandwidth}.
\end{example}

\begin{keyidea}[A model can be simple, but not simpler than the bandwidth you ask of it]
Every estimator trades speed against noise, and every sampled loop has a speed it cannot exceed. The useful bandwidth of an observer lies between the slowest disturbances it must catch and the fastest thing it can afford to believe: the noise of its sensor, the rate of its computer, the dynamics its model leaves out. Here the noise sets the practical limit near $\omega_o = 10$, the sampling a hard one near 160, and the motor lag, surprisingly, none.
\end{keyidea}

\screen[0 0 495 998]{adrcbw}{Lesson II.6 in the app with the observer at \SI{40}{rad/s}: the thrust in the motor chart is a band of noise, and the motors touch their limits.}{fig:adrcbw-screen}

\begin{inapp}
\begin{itemize}[leftmargin=1.2em]
  \item The lesson sets \param{control.adrc.wo = 40} and \param{sensors.posNoise = 0.02}; the wind is the default.
  \item The goal in the app computes the score over the last \SI{15}{s} from its own telemetry, sampled every \SI{5}{ms}; its jitter is therefore a little smaller than this chapter's, measured over \SI{10}{ms}, and its threshold is set accordingly.
  \item The observer gains are recomputed whenever $\omega_o$ changes; the discretisation with the matrix exponential (\src{src/estimation/eso.ts}) keeps the observer itself stable for any $\omega_o$ — it is the loop that becomes unstable.
\end{itemize}
\end{inapp}

\begin{trythis}
\begin{enumerate}
  \item Start at $\omega_o = 40$ and look at the thrust. Go down: 25, 15, 10, 6. Then up: 60, 100.
  \item Find the best score. Then halve the altimeter noise to \SI{1}{cm} and search again.
  \item Switch the noise off and raise $\omega_o$ to 150, then 200. Watch a setpoint step.
\end{enumerate}
\predict{if the altimeter noise is halved, by what factor does the thrust jitter fall at every $\omega_o$? And does the best $\omega_o$ move up or down?}
\end{trythis}

\section{The engineer's view: bandwidth is a budget}

\subsection{The same trade, everywhere}
Every feedback loop and every estimator faces the trade of this lesson, because of one identity (Appendix~C): the sensitivity $S$ and the complementary sensitivity $T$ add up to one at every frequency. Where a loop rejects disturbances ($S$ small) it must pass the sensor's signal faithfully ($T \approx 1$) — noise included. The bandwidth of a loop is the frequency up to which it believes its sensor. The Kalman filter of \lessonref{les:kalman} computed that frequency from the noise figures; here it is chosen by hand, and the score shows what the choice costs.

\subsection{Three limits on a bandwidth}
\begin{itemize}
  \item \textbf{Sensor noise} grows into the actuators as a power of the bandwidth — $\omega_o^{5/2}$ for an observer of a double integrator.
  \item \textbf{Sampling and delay} put a hard ceiling where the loop must react within a fraction of a sample (\lessonref{les:slow}).
  \item \textbf{Unmodelled dynamics} — a lag, a bending mode, a delay in the actuator — limit a loop that relies on a model of them. ADRC escapes the motor lag here because it lumps it into $f$; it would not escape a structural resonance, which is not slow.
\end{itemize}
A design is good when its bandwidth sits well below every one of them and well above the disturbances it must reject. When there is no such room, no tuning helps: a better sensor, a faster computer or a stiffer actuator does.

\begin{example}[The head of a hard-disk drive]
\label{ex:adrcbw-hdd}
The read head of a hard disk must stay over a track a few tens of nanometres wide while the disk spins. Its position is measured from servo marks written on the disk, read about \num{40000} times per second, and a disturbance observer cancels the forces of the air flow and of vibration. How fast can its observer be?

\textbf{Step 1 — the sampling limit.} With $\Delta t = \SI{25}{\micro s}$ and the edge of Example~\ref{ex:adrcbw-edge}, $\omega_o\Delta t < 0.65$: $\omega_o < \SI{26000}{rad/s}$, about \SI{4}{kHz}.

\textbf{Step 2 — the noise limit.} The position signal has noise of a few nanometres, a sizeable fraction of the tolerance. By \cref{eq:adrcbw-var}, every doubling of the bandwidth multiplies its effect on the actuator by 5.7.

\textbf{Step 3 — what is done.} Disk drives use loop bandwidths of one to two kilohertz — a factor of two to four below the sampling ceiling — and spend their engineering effort on the sensor and on knowing their resonances. The numbers differ from the drone's by a factor of a hundred; the reasoning is the same.
\end{example}

\subsection{In formal terms}
\Cref{thm:adrcbw-scaling} is one case of a general fact: for an estimator that recovers the $k$-th derivative of a measured signal with a bandwidth $\omega$, white measurement noise produces an error whose variance grows like $\omega^{2k+1}$. The observer of $f$ recovers the second derivative, hence $\omega^5$. The D term of Part~I recovers the first, hence $\omega^3$ — and \cref{eq:noise-simple}, $\sigma_D \propto f_c$, was the same law with the gain $K_d$ outside and the sample time fixed.

\begin{theorem}[Noise in an estimated derivative]
\label{thm:adrcbw-derivative}
Let $G(s) = \omega^{k+1}s^k/(s + \omega)^{k+1}$, an estimator of the $k$-th derivative with bandwidth $\omega$, be driven by white noise of intensity $q$. Then the output variance is $c_k\,q\,\omega^{2k+1}$ with $c_k = \frac{1}{2\pi}\int_{-\infty}^\infty \frac{x^{2k}}{(1 + x^2)^{k+1}}\,\dd x$, finite for every $k \ge 0$.
\end{theorem}
\begin{proof}
As for \cref{thm:adrcbw-scaling}, with the substitution $\omega' = \omega x$. The integrand decays like $x^{-2}$, so the integral converges.
\end{proof}

\begin{lookahead}
\begin{itemize}[leftmargin=1.2em]
  \item The accelerometer measures the second derivative directly, without amplifying the altimeter's noise: incremental dynamic inversion (Lessons~II.7 and II.8) gets the disturbance from a sensor instead of from an observer, and its noise problem becomes a filter problem (II.9).
  \item The trade between $S$ and $T$ is made quantitative by the sensitivity integral and loop shaping (Lessons~III.5 and III.6), and the observer's effect on the margins by Lesson~III.10.
  \item Quantisation is noise too: Lesson~III.19 sends a coarse sensor through the same derivative and finds a limit cycle; III.20 shows what aliasing adds.
\end{itemize}
\end{lookahead}

\begin{remember}
\begin{itemize}
  \item An observer cannot tell a disturbance from sensor noise: a faster one catches both.
  \item The noise in $\hat f$ grows like $\omega_o^{5/2}$; doubling the bandwidth multiplies it by 5.7.
  \item The thrust jitter is the altimeter noise through a linear filter; the sampled loop predicts it to two digits.
  \item With \SI{2}{cm} of noise the best $\omega_o$ is about 6; the hard ceiling, from sampling, is $\omega_o\Delta t \approx 0.65$ — here \SI{161}{rad/s}.
  \item Bandwidth is a budget between the disturbances to reject and the noise, sampling and unmodelled dynamics to ignore.
\end{itemize}
\end{remember}

\begin{curiosity}{Headphones that cancel the hum}
Noise-cancelling headphones are disturbance observers that you wear. A microphone inside each cup hears what reaches the ear; a feedback loop drives the loudspeaker with the opposite, so that the two cancel. On a plane the result is uncanny: the drone of the engines disappears, while voices and the clatter of the trolley come through almost unchanged.

The reason is this lesson's trade. The loop is fast — tens of thousands of samples per second — but the sound needs time to travel from the speaker to the microphone, a fraction of a millisecond, and that delay sets a ceiling on the bandwidth exactly as the sampling did for the drone. Below about a kilohertz the loop is far below its ceiling and cancels well; the low, steady hum of engines and air conditioning lives there. Above, the loop must give up, and the cups' passive padding takes over.

Push the bandwidth higher and the loop would start to amplify instead of cancel — the sensitivity bulge of \lessonref{les:challenge}. Engineers tune these loops so that the bulge falls where human hearing, and the music, care least.
\end{curiosity}

\begin{exercises}
\exercise[1] By \cref{eq:adrcbw-var}, by what factor does the noise in $\hat f$ change when $\omega_o$ goes from 10 to 15? Compare with the simulator's jitter, 0.371 and \SI{0.923}{N}.
\answer{$1.5^{5/2} = 2.76$; the simulator: $0.923/0.371 = 2.49$.}
\exercise[1] The altimeter noise is halved. By what factor does the thrust jitter fall at every $\omega_o$?
\answer{By two: the jitter is proportional to the noise's standard deviation.}
\exercise[1] What is the sampling ceiling of the observer if the controller runs at \SI{500}{Hz}? At \SI{100}{Hz}?
\answer{$0.65/0.002 = \SI{325}{rad/s}$; $0.65/0.01 = \SI{65}{rad/s}$.}
\exercise[2]\exkind{derive} Derive \cref{eq:adrcbw-G}: write the observer with the measurement $y + n$, take Laplace transforms, and solve for $\hat f$ in terms of $n$ (with $f = 0$, $u = 0$).
\answer{From the three rows: $s\hat y = \hat v + l_1(n - \hat y)$, $s\hat v = \hat f + l_2(n - \hat y)$, $s\hat f = l_3(n - \hat y)$. Eliminating: $(n - \hat y)(s^3 + l_1s^2 + l_2s + l_3) = s^3n$, so $\hat f = l_3s^2n/(s^3 + l_1s^2 + l_2s + l_3) = \omega_o^3s^2n/(s + \omega_o)^3$.}
\exercise[2]\exkind{derive} Compute $c_1$ of \cref{thm:adrcbw-derivative} and use it to estimate the noise of a velocity estimate with bandwidth \SI{20}{rad/s} from the altimeter of this lesson ($q = \sigma^2\Delta t$ with $\sigma = \SI{2}{cm}$, $\Delta t = \SI{4}{ms}$).
\answer{$\int x^2/(1 + x^2)^2\,\dd x = \pi/2$, so $c_1 = 1/4$. $\sigma_v^2 = \tfrac14 \cdot 1.6 \times 10^{-6} \cdot 20^3 = 3.2 \times 10^{-3}$: $\sigma_v = \SI{0.057}{m/s}$.}
\exercise[2]\exkind{derive} The best bandwidth in a model. Suppose the disturbance part of the score falls like $a/\omega_o$ and the noise part rises like $b\,\omega_o^{5/2}$. Find the minimising $\omega_o$, and how it moves when the noise is halved.
\answer{Derivative $-a/\omega_o^2 + \tfrac52b\,\omega_o^{3/2} = 0$: $\omega_o^* = (2a/5b)^{2/7}$. Halving the noise halves $b$: $\omega_o^*$ grows by $2^{2/7} = 1.22$.}
\exercise[2]\exkind{fly} In Lesson~II.6 switch the noise off and the wind on, and compare $\omega_o = 6$ and 30. Then switch the wind off and the noise on, and compare again. Which half of the score does each experiment isolate?
\answer{Wind without noise: only the disturbance half; the larger $\omega_o$ wins clearly (smaller error, no jitter). Noise without wind: only the noise half; the smaller $\omega_o$ wins (far less jitter, and the error is the small noise-induced motion).}
\exercise[2]\exkind{code} \src{src/estimation/eso.ts} discretises the observer with \texttt{c2d}, the exact zero-order-hold formula. What would happen at $\omega_o = 100$ with a forward-Euler discretisation instead? (Appendix~D, Example~\ref{ex:D-methods}.)
\answer{Forward Euler is stable only if every observer pole satisfies $|1 + s\Delta t| < 1$. The triple pole at $-100$ gives $1 - 0.4 = 0.6$: stable, but at $\omega_o > 500$ ($s\Delta t < -2$) the observer itself would diverge; with the exact discretisation it never does.}
\exercise[3]\exkind{derive} Why the motor lag does not matter here. The plant is $\ddot y = T$ with $\tau_m\dot T + T = u$, and the observer assumes $\ddot y = u + f$. Show that the lag appears in $f$ as $f = T - u = -\tau_m\dot T$, and argue why an observer much faster than $1/\tau_m$ can still estimate it.
\answer{$\ddot y = T = u + (T - u)$, and $T - u = -\tau_m\dot T$ by the lag's equation. This $f$ is a filtered derivative of the thrust, which changes on the time scale $\tau_m$; an observer with $\omega_o \gg 1/\tau_m$ tracks it with small error, and cancelling it simply makes the loop ignore the lag. What it cannot track is what changes faster than the sampling allows.}
\end{exercises}
